\begin{table}[htbp]
\centering
\small
\caption{V2 H1--H2（幾何再編と表現共有度の変位）：事後学習に伴う表現幾何学的歪み、デコードピーク深度変位 $\Delta d^*$、およびタスク共有度変化 $\Delta\text{Sharing}$。4ファミリー統合ブートストラップ95\%信頼区間。}
\label{tab:v2_h1_h2_reorganization}
\begin{tabular}{ll ccc}
\toprule
\textbf{Hypothesis} & \textbf{Metric} & \textbf{Estimate} & \textbf{95\% CI} & \textbf{Condition} \\
\midrule
\multirow{4}{*}{H1a: Geometric Distortion} & Reader Distortion ($1 - \text{CKA}$)     & 0.352    & [0.330, 0.370]     & Matched-Plain \\
                                 & Self Distortion ($1 - \text{CKA}$)       & 0.425    & [0.402, 0.445]     & Matched-Plain \\
                                 & RSA Reader ($\rho_{\text{RSA}}$)         & 0.648    & [0.630, 0.670]     & Matched-Plain \\
                                 & RSA Self ($\rho_{\text{RSA}}$)           & 0.575    & [0.555, 0.598]     & Matched-Plain \\
\midrule
\multirow{4}{*}{H1b: Decodability Peak Shift} & Valence Reader $\Delta d^*$              & 0.112    & [0.090, 0.138]     & Matched-Plain \\
                                 & Valence Self $\Delta d^*$                & 0.077    & [0.060, 0.100]     & Matched-Plain \\
                                 & Arousal Reader $\Delta d^*$              & 0.090    & [0.070, 0.110]     & Matched-Plain \\
                                 & Arousal Self $\Delta d^*$                & 0.060    & [0.045, 0.080]     & Matched-Plain \\
\midrule
\multirow{2}{*}{H2: Sharing Reorganization} & $\Delta\text{Sharing}$ (Valence)         & ---      & [-0.105, -0.060]   & Matched-Plain \\
                                 & $\Delta\text{Sharing}$ (Arousal)         & ---      & [-0.065, -0.045]   & Matched-Plain \\
\bottomrule
\end{tabular}
\vspace{1ex}
\begin{minipage}{\linewidth}
\footnotesize
\textbf{Note:} H1bにおいてすべての信頼区間下限が正（$\text{CI}_{\text{low}} > 0$）であり、事後学習によって感情表現のデコードピークが有意に浅層側（入力層側）へシフトした。またH2において $\Delta\text{Sharing} < 0$ であり、事後学習に伴い他者認識と自己報告の表現共有度がわずかに分離（特化）した。
\end{minipage}
\end{table}